%! Absolute Value Equations \& Inequalities — Unit 2, Lesson 2.2 — Homework
\documentclass[10pt]{article}
\usepackage{algebra2-article}
\usepackage{algebra2-boxes}
\newcommand{\TallMath}[1]{$\displaystyle #1\rule[-1.4em]{0pt}{3.2em}$}
\newcommand{\numline}[2]{%
  \draw[->, semithick, charcoal] (#1-0.4,0)--(#2+0.4,0);
  \foreach \x in {#1,...,#2}{\draw[charcoal] (\x,0.10)--(\x,-0.10) node[below, font=\scriptsize]{$\x$};}}

\begin{document}

\pageheader{Unit 2, Lesson 2.2}{Homework}
\namedateperiod

\vspace{0.10in}

\begin{notesbox}{Practice}
Show your work. Isolate the absolute value \emph{before} splitting, and write every inequality
solution in set, interval, and number-line form.

\begin{enumerate}[leftmargin=1.9em, itemsep=11pt]
  \item \textbf{Solve (already isolated).}
    \begin{enumerate}[label=(\alph*), leftmargin=1.8em, itemsep=4pt]
      \item $|x| = 10$: \quad $x = \blank{3.0cm}$
      \item $|x - 5| = 2$: \quad $x = \blank{3.0cm}$
    \end{enumerate}

  \item \textbf{Isolate, then solve.}
    \begin{enumerate}[label=(\alph*), leftmargin=1.8em, itemsep=4pt]
      \item $|4x + 2| = 10$: \quad $x = \blank{3.0cm}$
      \item $2|x - 4| - 1 = 9$: \quad $x = \blank{3.0cm}$
    \end{enumerate}

  \item \textbf{How many solutions?} Write \emph{none}, \emph{one}, or \emph{two}, and give any
        solution(s).
    \begin{enumerate}[label=(\alph*), leftmargin=1.8em, itemsep=4pt]
      \item $|x + 7| = -2$: \blank{4.5cm}
      \item $|x - 3| = 0$: \blank{4.5cm}
      \item $|x| = 5$: \blank{4.5cm}
    \end{enumerate}

  \item \textbf{Solve the ``and'' inequality} $|x + 2| < 6$.
        \\[4pt] Solve: \blank{4.5cm} \qquad Set: \blank{3.4cm} \qquad Interval: \blank{2.8cm}
        \begin{center}
        \begin{tikzpicture}[scale=0.5]
          \numline{-9}{5}
        \end{tikzpicture}
        \end{center}

  \item \textbf{Solve the ``or'' inequality} $|2x - 1| \ge 7$.
        \\[4pt] Solve: \blank{4.5cm} \qquad Set: \blank{3.4cm} \qquad Interval: \blank{3.2cm}
        \begin{center}
        \begin{tikzpicture}[scale=0.5]
          \numline{-5}{6}
        \end{tikzpicture}
        \end{center}

  \item \textbf{Model it.} A thermostat holds a room within $2^\circ$F of the $68^\circ$F setting.
    \begin{enumerate}[label=(\alph*), leftmargin=1.8em, itemsep=4pt]
      \item Write an absolute-value inequality for an acceptable temperature $T$. \blank{4.0cm}
      \item Solve it and state the coolest and warmest acceptable temperatures. \blank{5.0cm}
    \end{enumerate}

  \item \textbf{Explain.} In your own words, how do you decide whether an absolute-value inequality
        becomes an ``and'' (one interval) or an ``or'' (two rays)? \writelines{2}
\end{enumerate}
\end{notesbox}

\begin{extensionbox}
\textbf{Tolerance and reasoning.}
\begin{enumerate}[label=(\alph*), leftmargin=1.9em, itemsep=6pt]
  \item A bolt is acceptable if its diameter $d$ is within $0.2$ mm of the $15$ mm target. Write and
        solve an inequality; give the smallest and largest acceptable diameters.
  \item For what values of $k$ does $|x - 3| < k$ have \emph{no solution}? For what values of $k$ is
        \emph{every} real number a solution of $|x - 3| > k$? Justify from the meaning of absolute
        value.
\end{enumerate}
\writelines{3}
\end{extensionbox}

\begin{spiralbox}
\textbf{Coming next --- Lesson 2.3: Absolute Value Functions \& Transformations.} Today you
\emph{solved} $|ax+b|=c$; next you'll \emph{graph} $y=|ax+b|$ --- a \textbf{V-shape}. The two
solutions you found for an equation are exactly where that V crosses the horizontal line $y=c$, which
previews the vertex, axis of symmetry, and minimum you'll study formally in Unit~3.
\end{spiralbox}

\end{document}
